% -*- mode: latex; fill-column: 120; -*-
\OpenCHAMI{} is very configurable, and it has a client that allows you do most of the configuration
via the command line, so installing that client makes management much easier.
% begin_ohpc_run
% ohpc_comment_header OpenCHAMI Client Install \ref{sec:openchami_client_install}
\begin{lstlisting}[language=bash,keywords={}]
[sms](*\#*) dnf install -y \
	https://github.com/OpenCHAMI/ochami/releases/download/v0.3.4/ochami_0.3.4_linux_amd64.rpm
[sms](*\#*) yes | ochami config cluster set --system --default \
	${compute_prefix} cluster.uri https://${sms_name}:8443
\end{lstlisting}
% ohpc_validation_newline
% end_ohpc_run
